\documentclass[12pt,a4paper]{article}

\usepackage[utf8]{inputenc}
\usepackage[T1]{fontenc}
\usepackage{amsmath}
\usepackage{amssymb}
\usepackage{amsthm}
\usepackage{geometry}
\usepackage{enumitem}
\usepackage{booktabs}
\usepackage{array}
\usepackage{hyperref}
\usepackage{tikz}
\usepackage{float}

\geometry{margin=1in}

\newtheorem{definition}{Definition}
\newtheorem{theorem}{Theorem}

\begin{document}

%=========================================================
% TITLE PAGE
%=========================================================

\begin{titlepage}
\centering

```
\vspace*{5cm}

{\Huge\bfseries Ford--Fulkerson and Edmonds--Karp Algorithms\par}

\vspace{1.5cm}

{\Large\bfseries Design and Analysis of Algorithms (DAA)\par}

\vspace{1cm}

{\large\bfseries Name: Advaith Babu K\par}

\vspace{0.4cm}

{\large\bfseries Register Number: 24100943\par}

\vfill
```

\end{titlepage}

%=========================================================
% CONTENT
%=========================================================

\section{Introduction}

Many real-world problems can be represented using a network of connected
points. Examples include transportation networks, computer networks,
water pipelines, communication systems, and traffic networks.

In such problems, we may want to find the maximum amount of flow that
can be sent from a starting point to a destination point. The starting
point is called the source and the destination point is called the sink.

The Maximum Flow problem is one of the important problems in graph
algorithms. Two well-known algorithms used to solve this problem are:

\begin{enumerate}
\item Ford--Fulkerson Algorithm
\item Edmonds--Karp Algorithm
\end{enumerate}

Both algorithms repeatedly find an augmenting path and increase the
flow until no more possible paths are available.

\section{Flow Network}

A flow network is a directed graph in which every edge has a capacity.

A flow network can be represented as:

[
G=(V,E)
]

where:

\begin{itemize}
\item $V$ is the set of vertices.
\item $E$ is the set of directed edges.
\item $s$ is the source vertex.
\item $t$ is the sink vertex.
\item $c(u,v)$ is the capacity of edge $(u,v)$.
\item $f(u,v)$ is the current flow through edge $(u,v)$.
\end{itemize}

The flow must satisfy two important conditions.

\subsection{Capacity Constraint}

For every edge $(u,v)$,

[
0 \leq f(u,v) \leq c(u,v)
]

This means that the flow through an edge cannot be greater than its
capacity.

\subsection{Flow Conservation}

For every vertex other than the source and sink, the amount of incoming
flow must be equal to the amount of outgoing flow.

[
\sum\_{u} f(u,v)=\sum\_{w} f(v,w)
]

Therefore, flow cannot be created or destroyed at an intermediate
vertex.

\section{Residual Graph}

The residual graph is an important concept used by the
Ford--Fulkerson algorithm.

For every edge $(u,v)$, the residual capacity is:

[
c_f(u,v)=c(u,v)-f(u,v)
]

The residual graph also contains reverse edges. These reverse edges
allow the algorithm to cancel some previously assigned flow if a
better path is found later.

For example, if an edge has capacity $10$ and currently carries flow
$6$, then its remaining forward capacity is:

[
10-6=4
]

The corresponding reverse edge can have residual capacity $6$.

This allows the algorithm to send flow backwards and modify an earlier
decision.

\section{Augmenting Path}

An augmenting path is a path from the source $s$ to the sink $t$ in the
residual graph.

If an augmenting path exists, additional flow can be sent through that
path.

The amount of additional flow that can be sent is called the
bottleneck capacity.

It is given by:

[
\Delta =
\min\_{(u,v)\in P} c_f(u,v)
]

where $P$ is the augmenting path.

After finding the bottleneck, the flow along every edge of the path is
increased by $\Delta$.

\section{Ford--Fulkerson Algorithm}

The Ford--Fulkerson algorithm is a method for finding the maximum flow
in a flow network.

The main idea is simple:

\begin{enumerate}
\item Start with zero flow.
\item Find any augmenting path from the source to the sink.
\item Find the minimum residual capacity on that path.
\item Increase the flow along the path by that amount.
\item Update the residual graph.
\item Repeat until there is no augmenting path.
\end{enumerate}

When no augmenting path exists, the current flow is a maximum flow.

\subsection{Ford--Fulkerson Algorithm Steps}

\begin{enumerate}[label=\textbf{Step \arabic\*:}]

```
\item Initialize the flow on every edge to zero.

\item Construct the residual graph using the capacities of the
edges.

\item Find any path from the source $s$ to the sink $t$ in the
residual graph.

\item Find the minimum residual capacity on this path:

\[
\Delta =
\min_{(u,v)\in P} c_f(u,v)
\]

\item Increase the flow along the path by $\Delta$.

\item Update the residual capacities, including the reverse edges.

\item Repeat the process until there is no path from $s$ to $t$ in
the residual graph.

\item The resulting flow is the maximum flow.
```

\end{enumerate}

\section{Example of Ford--Fulkerson}

Consider the following flow network.

\begin{center}

\begin{tikzpicture}[
scale=1.2,
every node/.style={
circle,
draw,
minimum size=9mm
},
\>=stealth
]

\node (s) at (0,0) {$s$};
\node (a) at (2,1.5) {$A$};
\node (b) at (2,-1.5) {$B$};
\node (t) at (4,0) {$t$};

\draw[->] (s) -- node[above left] {$10$} (a);
\draw[->] (s) -- node[below left] {$10$} (b);
\draw[->] (a) -- node[above right] {$4$} (t);
\draw[->] (b) -- node[below right] {$6$} (t);
\draw[->] (a) -- node[right] {$8$} (b);

\end{tikzpicture}

\end{center}

Suppose the first augmenting path selected is:

[
s\rightarrow A\rightarrow t
]

The capacities on this path are:

[
10,\quad 4
]

Therefore, the bottleneck capacity is:

[
\Delta=\min(10,4)=4
]

The flow of $4$ units is sent through this path.

Next, suppose the algorithm selects:

[
s\rightarrow B\rightarrow t
]

The capacities available on this path are:

[
10,\quad 6
]

Therefore:

[
\Delta=\min(10,6)=6
]

The flow is increased by $6$ units.

The total flow is now:

[
4+6=10
]

The algorithm continues searching for augmenting paths until none are
available.

\section{Time Complexity of Ford--Fulkerson}

The running time of Ford--Fulkerson depends on how augmenting paths are
selected and on the capacities of the edges.

If all edge capacities are integers, every augmentation increases the
flow by at least one unit.

Let:

[
F=\text{value of the maximum flow}
]

Then there can be at most $F$ augmentations.

Finding an augmenting path using DFS or BFS takes:

[
O(E)
]

time.

Therefore, the total time complexity is:

[
\boxed{O(EF)}
]

for integer capacities.

\subsection{Important Point}

The complexity $O(EF)$ depends on the value of the maximum flow.
Therefore, Ford--Fulkerson is not considered a polynomial-time
algorithm when the capacities are represented in binary.

Also, with irrational capacities, Ford--Fulkerson can fail to terminate
for some choices of augmenting paths.

\section{Edmonds--Karp Algorithm}

Edmonds--Karp is a specific implementation of the
Ford--Fulkerson method.

The main difference is the way an augmenting path is selected.

Ford--Fulkerson can use any augmenting path.

Edmonds--Karp always selects the augmenting path with the smallest
number of edges. This path is found using Breadth-First Search (BFS).

Therefore:

[
\boxed{\text{Edmonds--Karp = Ford--Fulkerson + BFS}}
]

\section{Why BFS is Used}

BFS explores a graph level by level.

Therefore, when BFS reaches the sink, the path found contains the
minimum possible number of edges.

This gives Edmonds--Karp a predictable polynomial running time.

The algorithm does not depend on the numerical values of the
capacities for its running time.

\section{Edmonds--Karp Algorithm Steps}

\begin{enumerate}[label=\textbf{Step \arabic\*:}]

```
\item Initialize the flow on every edge to zero.

\item Construct the residual graph.

\item Use BFS to find the shortest augmenting path from the source
$s$ to the sink $t$.

\item If no path is found, stop. The current flow is the maximum
flow.

\item Find the bottleneck capacity of the path:

\[
\Delta =
\min_{(u,v)\in P} c_f(u,v)
\]

\item Increase the flow along the selected path by $\Delta$.

\item Update the residual capacities and reverse edges.

\item Repeat the BFS search until no augmenting path exists.

\item Return the resulting maximum flow.
```

\end{enumerate}

\section{Time Complexity of Edmonds--Karp}

In Edmonds--Karp, BFS is used to find every augmenting path.

The time required for one BFS is:

[
O(E)
]

A key property of Edmonds--Karp is that the number of augmentations is
bounded by:

[
O(VE)
]

Therefore:

[
O(E)\times O(VE)
]

gives:

[
\boxed{O(VE^2)}
]

Thus, the time complexity of Edmonds--Karp is:

[
\boxed{O(VE^2)}
]

where:

\begin{itemize}
\item $V$ is the number of vertices.
\item $E$ is the number of edges.
\end{itemize}

This is a polynomial-time complexity.

\section{Correctness of the Algorithms}

The correctness of both algorithms is based on the
Max-Flow Min-Cut theorem.

The theorem states that the maximum possible flow from the source to
the sink is equal to the minimum capacity of an $s$-$t$ cut.

\begin{theorem}
If there is no augmenting path from $s$ to $t$ in the residual graph,
then the current flow is a maximum flow.
\end{theorem}

\subsection{Reason}

The algorithm continuously increases the flow as long as an
augmenting path exists.

When no augmenting path remains, the vertices reachable from the
source in the residual graph form a cut.

All edges crossing this cut from the reachable side to the
unreachable side are fully saturated.

Therefore, the capacity of this cut is equal to the value of the
current flow.

By the Max-Flow Min-Cut theorem, the current flow must be a maximum
flow.

\section{Comparison of Ford--Fulkerson and Edmonds--Karp}

\begin{center}

\renewcommand{\arraystretch}{1.5}

\begin{tabular}{>{\raggedright\arraybackslash}p{4cm}
\>{\raggedright\arraybackslash}p{5cm}
\>{\raggedright\arraybackslash}p{5cm}}

\toprule

\textbf{Feature}
&
\textbf{Ford--Fulkerson}
&
\textbf{Edmonds--Karp}
\\

\midrule

Path selection
&
Any augmenting path
&
Shortest augmenting path in terms of edges
\\

Search method
&
DFS or any suitable method
&
BFS
\\

Running time
&
$O(EF)$ for integer capacities
&
$O(VE^2)$
\\

Dependence on capacities
&
Depends on maximum flow value
&
Independent of capacity values for the asymptotic bound
\\

Polynomial time
&
Not guaranteed in general
&
Yes
\\

Implementation
&
Simpler
&
More structured because BFS is used
\\

\bottomrule

\end{tabular}

\end{center}

\section{Important Terms}

\begin{description}

```
\item[Flow:] The amount of a commodity sent through a network.

\item[Capacity:] The maximum amount of flow that an edge can carry.

\item[Source:] The vertex from which the flow originates.

\item[Sink:] The vertex where the flow terminates.

\item[Residual Capacity:] The remaining capacity available on an
edge.

\item[Residual Graph:] A graph representing the remaining possible
flow and reverse flow.

\item[Augmenting Path:] A path from the source to the sink in the
residual graph.

\item[Bottleneck Capacity:] The smallest residual capacity on an
augmenting path.

\item[Maximum Flow:] The largest possible flow from the source to
the sink.
```

\end{description}

\section{Complexity Summary}

The main complexities can be summarized as follows.

\begin{center}

\renewcommand{\arraystretch}{1.5}

\begin{tabular}{>{\raggedright\arraybackslash}p{5cm}
\>{\raggedright\arraybackslash}p{5cm}}

\toprule

\textbf{Algorithm} & \textbf{Time Complexity} \\

\midrule

Ford--Fulkerson &
$O(EF)$ for integer capacities
\\

Edmonds--Karp &
$O(VE^2)$
\\

\bottomrule

\end{tabular}

\end{center}

Here, $V$ represents the number of vertices, $E$ represents the number
of edges, and $F$ represents the value of the maximum flow.

\section{Key Points}

The following points are important to remember:

\begin{enumerate}

```
\item Ford--Fulkerson is a method for finding maximum flow.

\item It starts with zero flow.

\item It repeatedly finds an augmenting path.

\item The bottleneck capacity determines how much additional flow
can be sent.

\item The residual graph contains both forward and reverse
capacities.

\item Ford--Fulkerson can choose any augmenting path.

\item For integer capacities, its complexity is:

\[
O(EF)
\]

\item Edmonds--Karp is a specific implementation of
Ford--Fulkerson.

\item Edmonds--Karp uses BFS to find the shortest augmenting path.

\item The time complexity of Edmonds--Karp is:

\[
O(VE^2)
\]

\item Edmonds--Karp always terminates in polynomial time.

\item Both algorithms are based on repeatedly increasing the flow
using augmenting paths.
```

\end{enumerate}

\section{Conclusion}

Ford--Fulkerson and Edmonds--Karp are important algorithms for solving
the Maximum Flow problem.

Ford--Fulkerson provides the basic idea of repeatedly finding
augmenting paths in a residual graph. Its performance depends on the
choice of paths and, for integer capacities, can be expressed as
$O(EF)$.

Edmonds--Karp improves the path-selection method by always using BFS to
find the shortest augmenting path. This gives it a polynomial
time complexity of:

[
O(VE^2)
]

The main difference to remember is that Ford--Fulkerson allows any
augmenting path, whereas Edmonds--Karp specifically uses BFS.

\end{document}

Close